% GENERATED FILE. Edit canonical lessons, then run npm run generate:books.
% canonical-source-hash: fnv1a64:e319dd9a9fed7b61
% canonical-lessons: TA-C266-elitaka, TA-C266-ventumenre, TA-C266-tarceyalaka, TA-C266-atirstavacamaka, TA-C266-karanam

\chapter{Easily, On purpose, By chance, Luckily, Reason}
\label{ch:ta-easily-on-purpose-by-chance-luckily-reason}
\hlchaptermodality{266}

\begin{chapteropening}
\textbf{By the end of this chapter:} \emph{I can say easily, on purpose, by chance, luckily and a reason.}

Complete the last lesson of chapter 266: I can say easily, on purpose, by chance, luckily and a reason.
\end{chapteropening}

\section[\texorpdfstring{\ta{எளிதாக} (eḷitāka)}{எளிதாக (eḷitāka)}]{\ta{எளிதாக} (eḷitāka) — easily}

\label{lesson:TA-C266-elitaka}



\begin{quote}
Before the new one: say the Tamil for a prison, then the Tamil for the world.
\end{quote}

\subsection*{You'll want to know: \ta{எளிதாக}}
\textbf{\ta{எளிதாக}} — \emph{eḷitāka} — \textquotedblleft{}easily\textquotedblright{}.

\begin{grammarlens}[title={What you've built}]
One more word for this chapter.
\end{grammarlens}

\subsection*{Your turn}
\begin{itemize}
\raggedright
  \item \emph{Say it:} \emph{eḷitāka}
  \item \emph{Say it:} \emph{eḷitāka}, once more
  \item \emph{Say it:} say \emph{eḷitāka}
  \item \emph{Recall:} say the Tamil for a prison, then the Tamil for the world, then say \emph{eḷitāka} again
\end{itemize}

\begin{tcolorbox}[breakable,colback=teal!4,colframe=teal!35!black,title={Before you move on}]
What does \ta{எளிதாக} mean? (\textquotedblleft{}Easily\textquotedblright{}.) Say it once more.
\end{tcolorbox}
\section[\texorpdfstring{\ta{வேண்டுமென்றே} (vēṇṭumeṉṟē)}{வேண்டுமென்றே (vēṇṭumeṉṟē)}]{\ta{வேண்டுமென்றே} (vēṇṭumeṉṟē) — on purpose, deliberately}

\label{lesson:TA-C266-ventumenre}



\begin{quote}
Before the new one: say the Tamil for the world, then the Tamil for easily.
\end{quote}

\subsection*{You'll want to know: \ta{வேண்டுமென்றே}}
\textbf{\ta{வேண்டுமென்றே}} — \emph{vēṇṭumeṉṟē} — \textquotedblleft{}on purpose, deliberately\textquotedblright{}.

\begin{grammarlens}[title={What you've built}]
One more word for this chapter.
\end{grammarlens}

\subsection*{Your turn}
\begin{itemize}
\raggedright
  \item \emph{Say it:} \emph{vēṇṭumeṉṟē}
  \item \emph{Say it:} \emph{vēṇṭumeṉṟē}, once more
  \item \emph{Say it:} say \emph{vēṇṭumeṉṟē}
  \item \emph{Recall:} say the Tamil for the world, then the Tamil for easily, then say \emph{vēṇṭumeṉṟē} again
\end{itemize}

\begin{tcolorbox}[breakable,colback=teal!4,colframe=teal!35!black,title={Before you move on}]
What does \ta{வேண்டுமென்றே} mean? (\textquotedblleft{}On purpose, deliberately\textquotedblright{}.) Say it once more.
\end{tcolorbox}
\section[\texorpdfstring{\ta{தற்செயலாக} (taṟceyalāka)}{தற்செயலாக (taṟceyalāka)}]{\ta{தற்செயலாக} (taṟceyalāka) — by chance}

\label{lesson:TA-C266-tarceyalaka}



\begin{quote}
Before the new one: say the Tamil for easily, then the Tamil for on purpose.
\end{quote}

\subsection*{You'll want to know: \ta{தற்செயலாக}}
\textbf{\ta{தற்செயலாக}} — \emph{taṟceyalāka} — \textquotedblleft{}by chance\textquotedblright{}.

\begin{grammarlens}[title={What you've built}]
One more word for this chapter.
\end{grammarlens}

\subsection*{Your turn}
\begin{itemize}
\raggedright
  \item \emph{Say it:} \emph{taṟceyalāka}
  \item \emph{Say it:} \emph{taṟceyalāka}, once more
  \item \emph{Say it:} say \emph{taṟceyalāka}
  \item \emph{Recall:} say the Tamil for easily, then the Tamil for on purpose, then say \emph{taṟceyalāka} again
\end{itemize}

\begin{tcolorbox}[breakable,colback=teal!4,colframe=teal!35!black,title={Before you move on}]
What does \ta{தற்செயலாக} mean? (\textquotedblleft{}By chance\textquotedblright{}.) Say it once more.
\end{tcolorbox}
\section[\texorpdfstring{\ta{அதிர்ஷ்டவசமாக} (atirṣṭavacamāka)}{அதிர்ஷ்டவசமாக (atirṣṭavacamāka)}]{\ta{அதிர்ஷ்டவசமாக} (atirṣṭavacamāka) — luckily, fortunately}

\label{lesson:TA-C266-atirstavacamaka}



\begin{quote}
Before the new one: say the Tamil for on purpose, then the Tamil for by chance.
\end{quote}

\subsection*{You'll want to know: \ta{அதிர்ஷ்டவசமாக}}
\textbf{\ta{அதிர்ஷ்டவசமாக}} — \emph{atirṣṭavacamāka} — \textquotedblleft{}luckily, fortunately\textquotedblright{}.

\begin{grammarlens}[title={What you've built}]
One more word for this chapter.
\end{grammarlens}

\subsection*{Your turn}
\begin{itemize}
\raggedright
  \item \emph{Say it:} \emph{atirṣṭavacamāka}
  \item \emph{Say it:} \emph{atirṣṭavacamāka}, once more
  \item \emph{Say it:} say \emph{atirṣṭavacamāka}
  \item \emph{Recall:} say the Tamil for on purpose, then the Tamil for by chance, then say \emph{atirṣṭavacamāka} again
\end{itemize}

\begin{tcolorbox}[breakable,colback=teal!4,colframe=teal!35!black,title={Before you move on}]
What does \ta{அதிர்ஷ்டவசமாக} mean? (\textquotedblleft{}Luckily, fortunately\textquotedblright{}.) Say it once more.
\end{tcolorbox}
\section[\texorpdfstring{\ta{காரணம்} (kāraṇam)}{காரணம் (kāraṇam)}]{\ta{காரணம்} (kāraṇam) — a reason, a cause}

\label{lesson:TA-C266-karanam}



\begin{quote}
Before the new one: say the Tamil for by chance, then the Tamil for luckily.
\end{quote}

\subsection*{You'll want to know: \ta{காரணம்}}
\textbf{\ta{காரணம்}} — \emph{kāraṇam} — \textquotedblleft{}a reason, a cause\textquotedblright{}.

\begin{grammarlens}[title={What you've built}]
One more word for this chapter.
\end{grammarlens}

\subsection*{Your turn}
\begin{itemize}
\raggedright
  \item \emph{Say it:} \emph{kāraṇam}
  \item \emph{Say it:} \emph{kāraṇam}, once more
  \item \emph{Say it:} say \emph{kāraṇam}
  \item \emph{Recall:} say the Tamil for by chance, then the Tamil for luckily, then say \emph{kāraṇam} again
\end{itemize}

\begin{tcolorbox}[breakable,colback=teal!4,colframe=teal!35!black,title={Before you move on}]
What does \ta{காரணம்} mean? (\textquotedblleft{}A reason, a cause\textquotedblright{}.) Say it once more.
\end{tcolorbox}
